% Part:sets-functions-relations
% Chapter: size-of-sets
% Section: equinumerous-sets

\documentclass[../../../include/open-logic-section]{subfiles}

\begin{document}

\olfileid{sfr}{siz}{equ}

\olsection{Evenveel elementen}

Bij eindige verzamelingen werkt ons gewone idee van ``grootte'' prima. Maar
hoe zit het met oneindige verzamelingen? Als we een precieze regel willen om
verzamelingen van \emph{elke} grootte te vergelijken, is het handig om
eerst te zeggen wanneer twee verzamelingen even groot zijn. Frege geeft dit voorbeeld:
\begin{quote}
  Als een ober zeker wil weten dat hij precies evenveel messen als borden op
  tafel heeft gelegd, hoeft hij geen van beide te tellen. Hij legt gewoon rechts
  van elk bord een mes en zorgt dat elk mes rechts van een bord ligt. Zo hoort
  bij elk bord precies één mes en bij elk mes precies één bord: het mes ligt
  telkens rechts van het bord. \citep[\S70]{Frege1884}
\end{quote}
Dit idee kunnen we nu precies opschrijven:

\begin{defn}\ollabel{comparisondef}
  $A$ is \emph{gelijkmachtig} met $B$, geschreven als $\cardeq{A}{B}$,
  precies wanneer er !!a{bijection} $f \colon A \to B$ bestaat. Zo'n functie
  koppelt de elementen één-op-één aan elkaar en slaat niets over.
\end{defn}

\begin{prop}\ollabel{equinumerosityisequi}
De relatie ``evenveel elementen hebben'' is een equivalentierelatie.
\end{prop}

\begin{proof} 
We moeten drie dingen laten zien: gelijkmachtigheid is reflexief, symmetrisch
en transitief. Neem verzamelingen $A, B$ en $C$.

\emph{Reflexiviteit.} De functie $\Id{A} \colon A \to A$ stuurt elk element
naar zichzelf: $\Id{A} (x) = x$ voor elke $x \in A$. Deze functie is
!!a{bijection}. Dus
$\cardeq{A}{A}$.

\emph{Symmetrie.} Stel dat $\cardeq{A}{B}$. Er is dan !!a{bijection}
$f\colon A \to B$. Omdat $f$ !!{bijective} is, bestaat de omgekeerde functie
$f^{-1}$ en is die ook !!{bijective}. Daarom is
$f^{-1}\colon B \to A$ !!a{bijection}, en dus $\cardeq{B}{A}$.

\emph{Transitiviteit.} Stel dat $\cardeq{A}{B}$ en $\cardeq{B}{C}$. Dan zijn
er !!{bijection}s $f\colon A \to B$ en $g\colon B \to C$. Voeren we die na
elkaar uit, dan krijgen we de bijectieve functie $\comp{f}{g}\colon A \to C$.
Daarom geldt $\cardeq{A}{C}$.
\end{proof}

\begin{prop}
Als $\cardeq{A}{B}$, dan is $A$ precies dan !!{enumerable} als $B$ dat is.
\end{prop}

\begin{editorial}
Het volgende bewijs gebruikt \olref[enm]{defn:enumerable} wanneer
\olref[enm]{sec} is opgenomen. Anders gebruikt het
\olref[enm-alt]{defn:enumerable}.
\end{editorial}

\begin{proof}
Stel dat $\cardeq{A}{B}$. Dan is er !!a{bijection} $f \colon A
\to B$. Stel ook dat $A$ !!{enumerable} is.
\oliflabeldef{sfr:siz:enm:defn:enumerable}{
  Dan is $A = \emptyset$, of er is een surjectieve functie
  $g\colon \PosInt \to A$. Als $A = \emptyset$, dan is ook $B = \emptyset$
  (anders zou er !!a{element}~$y \in B$ zijn, terwijl er geen $x \in
  A$ is met $f(x) = y$). In het andere geval is $g\colon \PosInt \to A$
  !!{surjective}. Dan is $\comp{g}{f} \colon \PosInt \to B$ ook
  !!{surjective}. Neem om dat te zien een $y \in B$. Omdat $f$
  !!{surjective} is, is er een $x \in A$ met $f(x) = y$. Omdat
  $g$ !!{surjective} is, is er een $n \in \PosInt$ met $g(n) =
  x$. Dus
\[
(\comp{g}{f})(n) = f(g(n)) = f(x) = y
\]
en daarmee is $\comp{g}{f}$ !!{surjective}. De functie $\comp{g}{f}$ somt
dus de hele verzameling~$B$ op. Daarom is $B$~!!{enumerable}.}
{
  Dan is $A = \emptyset$, of er is !!a{bijection}~$g$ waarvan het bereik
  $A$ is en het domein $\Nat$ of een eerste stuk van de natuurlijke getallen. Als
  $A = \emptyset$, dan is ook $B = \emptyset$ (anders zou er een~$y \in B$
  zijn zonder een $x \in A$ waarvoor $f(x) = y$). Neem dus zo'n
  !!{bijection}~$g$. Dan is $\comp{g}{f}$ !!a{bijection} met bereik~$B$ en
  met hetzelfde domein als~$g$ (dus $\Nat$ of een eerste stuk daarvan).
  Daarom is $B$ !!{enumerable}.}

Als $B$ !!{enumerable} is, laat hetzelfde argument zien dat $A$
!!{enumerable} is. We gebruiken dan de !!{bijection} $f^{-1}\colon B \to A$
in plaats van~$f$.
\end{proof}

\begin{prob}
Laat zien dat als $\cardeq{A}{C}$ en $\cardeq{B}{D}$, en $A \cap B =
C \cap D = \emptyset$, dan $\cardeq{A \cup B}{C \cup D}$.
\end{prob}

\begin{prob}
Laat zien dat als $A$ oneindig en !!{enumerable} is, dan
$\cardeq{A}{\Nat}$.
\end{prob}
\end{document}
